\section{Introduction}
\label{sec:intro}

Microscopic traffic simulation is widely used to evaluate congestion, signal
control, routing, emissions, and network-level mobility.  Simulators such as
SUMO provide detailed measurements and allow interventions to be tested without
affecting a real transportation system.  However, simulation is primarily an
evaluation tool: converting a diagnosed bottleneck into an executable action
still requires an analyst to select a control strategy, modify interdependent
simulation artifacts, and evaluate possible network-wide side effects.  This
gap between diagnosis and action limits the automation of simulation-based
traffic optimization.

Large language models (LLMs) offer a flexible way to reason across heterogeneous
interventions, including signal retiming, corridor control, and demand
rerouting.  Nevertheless, plausible language is not evidence of traffic
benefit.  An LLM may reference nonexistent network entities, request unsupported
operations, omit required parameters, or propose an intervention that improves
a local bottleneck by transferring congestion elsewhere.  Traffic actions may
also interact: two individually reasonable interventions can reinforce or
interfere with one another.  Because exhaustive simulation of all combinations
is impractical, an LLM-based optimizer must ground proposals in executable
capabilities, search combinations within a limited budget, and reserve final
selection for measured simulation outcomes.

This study develops a simulation-grounded multi-agent LLM traffic optimizer.
The framework first simulates the unmodified scenario and extracts compact
bottleneck-centered evidence containing traffic measurements, local
connectivity, and relevant signal identifiers.  Three specialist agents then
generate structured proposals for signal control, corridor-speed modification,
and demand-preserving rerouting.  A deterministic capability registry rejects
unsupported actions, invalid identifiers, and incomplete parameters before
execution.  Every valid proposal is applied to an isolated scenario and
evaluated through a complete SUMO simulation.  The LLM therefore serves as a
candidate generator and local search operator rather than as the final
traffic-control decision maker.

The system further performs a simulation-feedback beam search.  Measured KPIs
and constraint violations are returned to the responsible specialist, which
produces an auditable one-parameter refinement of a retained parent.  The
remaining budget is used for pairwise search over compatible interventions and
deterministic enumeration of otherwise-unexplored local neighbors.
Candidate actions are applied sequentially and re-simulated, so their combined
effect is measured rather than inferred from textual predictions.  Final selection minimizes time loss on ten baseline-frozen
bottleneck edges while preventing increases in network-wide delay, excessive
spillover to neighboring roads, throughput loss, or additional teleport events.
The unmodified scenario remains available as a no-operation guard.  An offline
Pareto audit additionally checks whether the delay-selected plan is dominated
on measured delay, emissions, or mobility outcomes.

The primary evaluation covers five urban networks, three demand seeds, and
2,400 origin--destination trips per scenario.  Across 15 runs, the selected
outcomes reduce Top-10 bottleneck time loss by $9.18\%\pm10.18\%$ and
network-wide delay by 2.78\%.  The LLM-feedback search finds a feasible
improving plan in 86.7\% of cases, compared with 73.3\% for grounded random
search and 80.0\% for a deterministic heuristic under the same simulation cap.
Eleven of its 15 final selections are nondominated within their evaluated
candidate pools.

This work contributes a simulation-grounded multi-agent framework that generates
heterogeneous traffic interventions from compact diagnostic and topological
evidence while validating every proposal against executable capabilities.  It
introduces a budgeted composition mechanism that applies and re-simulates
compatible intervention pairs, allowing measured interactions to influence
selection without exhaustive search.  The framework further combines a
baseline-frozen bottleneck objective with constraints on global delay,
spillover, throughput, and simulation integrity, together with a no-op guard
and an offline Pareto dominance audit.  Finally, the multi-network, multi-seed
evaluation quantifies both achievable improvements and the rejection of locally
promising but globally unsafe candidates.  Together, these components position
the LLM as an adaptable proposal and reasoning module while reserving final
authority for executable tools, simulation evidence, and explicit optimization
constraints.
